\section{Contribution}

The main contribution of this thesis is a map of where AI can be used across the workflow around GGM, and of the reasoning behind each placement.
The workflow is followed step by step, from obtaining source data through preparing it, setting up a scenario and calibrating the model, to running it and reading the results.
At each step the thesis asks what happens today, what the step requires, what an agent could do, and where a person should stay involved.
The result places each step somewhere between work that ordinary code should do, work an agent could take on, and work that should remain with the analyst, and ends with a description of each proposed agent: what it receives, what it produces, what tools it uses, what it must never do, and how one would know it worked.
That description is meant to be the starting point for building the agents in the master's thesis.

The thesis also makes two contributions beyond GGM.
First, it carries a division of labour that the recent literature supports with evidence (language models reason and propose, deterministic code computes, people decide where decisions are hard to undo) from engineering and empirical economics to a market equilibrium model, and tests it against a real workflow step by step.
Second, it identifies the calibration of models whose parameters encode economic judgment as a case the existing literature on agent-based calibration does not cover.
In such a case, keeping a person in charge of the decision is not only a safety margin that better models will make unnecessary, but follows from what is being decided.

The contribution is also practical.
The work of preparing and calibrating GGM is a real and documented burden for the researchers who use it, and a clear account of which parts of it could be automated, and at what risk, is useful to them in its own right.
What the thesis does not contribute is evidence that any particular agent works; that is left to the master's thesis.
